% NS-55 prime-shift lane: exact identities and scoped operator obstructions.
% Uses v118 physical realization and v154 affine-potential notation.
\subsection*{Prime-shift primitives and the limits of transfer arguments}

\paragraph{Scope.}
Let $I_a=(-a,a)$, $L=2a$, and use the complete operator $A_a$,
form domain $V_a$ and screw function $\mathfrak g$ of
Proposition~\ref{prop:ns51-affine-potential}. Write
$\ell_n=\log n$ and $w_n=\Lambda(n)/\sqrt n$. Terms with
$w_n=0$ are omitted. The following are exact identities and proved
obstructions to particular inferences. They prove neither an arithmetic
nullvector nor affine-potential injectivity (API). Both pole signs are
retained. No new window or tail metric is computed; G2 and RH remain open.

\begin{proposition}[The exact two-sided primitive equation]
\label{prop:ns55-prime-primitives}
Remove the prime ramps from the screw function, writing on $|t|<L$
\[
 \mathfrak g(t)=\mathfrak g_0(t)+
       \sum_{1<n<e^L}w_n(|t|-\ell_n)_+.
\]
Thus $\mathfrak g_0$ includes the archimedean and both pole terms.
For $f\in L^2(I_a)$ define the clamped primitive and total integral
\[
 H_f(x)=\int_{-\infty}^x E_af(y)\,dy,
 \qquad M_f=\int_{I_a}f(y)\,dy.
\]
Then, at every $x\in I_a$,
\begin{equation}\label{eq:ns55-prime-primitive}
 U_f(x)=(\mathfrak g_0'*E_af)(x)+
 \sum_{1<n<e^L}w_n
       \{H_f(x-\ell_n)+H_f(x+\ell_n)-M_f\}.
\end{equation}
In particular API's equation $U_f=b$ is exactly this two-sided
shifted-primitive equation, subject to $f\in V_a$ and the parity constants
of Proposition~\ref{prop:ns51-affine-potential}. Differentiation is valid
in distributions and gives
\begin{equation}\label{eq:ns55-prime-differentiated}
 (\mathfrak g_0''*E_af)|_{I_a}
 +\sum_{1<n<e^L}w_n\{E_af(\cdot-\ell_n)+E_af(\cdot+\ell_n)\}|_{I_a}=0.
\end{equation}
No $H^1$ regularity of $f$ is required or inferred.
\end{proposition}
\begin{proof}
The derivative of $(|t|-\ell)_+$ is
$\mathbf1_{t>\ell}-\mathbf1_{t<-\ell}$ almost everywhere. Its
convolution with $E_af$ is
$H_f(x-\ell_n)-\{M_f-H_f(x+\ell_n)\}$.
The local $L^2$ property of $\mathfrak g_0'$ and compact support of
$E_af$ make its convolution continuous, as in the affine-potential
criterion. The finite sum is continuous too. Finally
$H_f'=E_af$ in distributions. Membership of this antiderivative in
$H^1_{\rm loc}$ does not imply $f\in H^1$.
\end{proof}

\paragraph{Why this is not a propagation rule.}
The prime part couples both $x-\ell_n$ and $x+\ell_n$.
More importantly, $\mathfrak g_0''*E_af$ remains a nonlocal term.
Knowing $f=0$ on a subinterval does not make that term zero there;
values of $f$ throughout the interval still enter. The logarithmic
singularity also lies on the convolution diagonal, so analyticity
of the kernel away from its singularities does not imply analyticity
of the potential across the support. No identity theorem can be
applied to that unsupported analytic continuation. Additional
prime-free exterior hypotheses can change this conclusion; none
is imposed in the present statement.

\begin{proposition}[The full semigroup is not positivity preserving]
\label{prop:ns55-nonpositive-semigroup}
Let $r>1$ be the unique real solution of $r^3-r-1=0$, and suppose
$L>\log r$. There are real nonnegative
$f,g\in C_c^\infty(I_a)$ with disjoint supports such that
\begin{equation}\label{eq:ns55-positive-cross}
 q_a(f,g)>0.
\end{equation}
Consequently the complete semigroup $e^{-sA_a}$ is not positivity
preserving on the full real $L^2(I_a)$ space with its ordinary
nonnegative cone: for these tests and all sufficiently small $s>0$,
\[
 \langle e^{-sA_a}f,g\rangle<0.
\]
This asserts no negative diagonal Weil energy and makes no separate
claim about order preservation after a parity restriction.
\end{proposition}
\begin{proof}
For disjoint supports the diagonal terms in the operator contribute
nothing. Away from prime-power separations, the off-diagonal integral
kernel of Proposition~\ref{prop:v118-log-dirichlet} is
\[
 J(t)=2\cosh(t/2)-\rho(t),\qquad
 \rho(t)=\frac{e^{t/2}}{e^t-e^{-t}},\quad t>0.
\]
The logarithmic Laplacian contributes $-1/(2t)$ and its bounded
correction contributes $-(\rho(t)-1/(2t))$, giving $-\rho(t)$;
the pole kernel gives the displayed $2\cosh(t/2)$.
For $u=e^t>1$,
\[
 J(t)=\frac{u^3-u-1}{\sqrt u\,(u^2-1)}.
\]
The numerator is strictly increasing on $(1,\infty)$, changes sign
once, and is positive precisely when $t>\log r$.
Choose a separation $t_0\in(\log r,L)$ outside the finite set
$\{\log n:1<n<e^L,\ \Lambda(n)\ne0\}$. Two sufficiently small
disjoint intervals inside $I_a$, at separation $t_0$, have all
cross separations in the positive region of $J$ and avoid every
prime-power separation. Choose nonzero nonnegative smooth bumps
$f,g$ supported in these intervals. All translated prime cross
terms vanish and the remaining double integral is strictly positive.
The tests belong to $\mathcal D(A_a)$, and $A_a$ is self-adjoint
and bounded below at this fixed window. Strong differentiability
of its semigroup on the operator domain gives
\[
 \langle e^{-sA_a}f,g\rangle
 =\langle f,g\rangle-s\langle A_af,g\rangle+o(s)
 =-s q_a(f,g)+o(s),
\]
which is negative for small positive $s$.
\end{proof}

\paragraph{Meaning of the order obstruction.}
A standard positive-kernel or maximum-principle transfer cannot simply
be attributed to the complete arithmetic operator. This does not
exclude a more specialized uniqueness argument, positivity of a
particular eigenfunction, a different cone, or API itself.
Positivity of the quadratic form is distinct from preservation of
pointwise positivity by its semigroup. Moreover conjugation by a
strictly positive multiplier $h$ with $h,h^{-1}\in L^\infty(I_a)$
cannot repair this: multiplication by $h$ and by $h^{-1}$ both
preserve the nonnegative cone, so positivity of the conjugated
semigroup would imply positivity of the original one. This concerns
the full signed operator, not a separate positive edge form.

\begin{proposition}[Prime entry is not an analytic window perturbation]
\label{prop:ns55-prime-entry}
Rescale $(0,L)$ unitarily to $(0,1)$. For a single prime-power length
$\ell>0$, define the symmetric compressed shift
\[
 (J_{\ell,L}u)(x)=
 \mathbf1_{x+\ell/L<1}u(x+\ell/L)
 +\mathbf1_{x>\ell/L}u(x-\ell/L).
\]
It vanishes for $L\le\ell$, but
\begin{equation}\label{eq:ns55-prime-entry-norm}
 \|J_{\ell,L}\|_{2\to2}=1\qquad(\ell<L<2\ell).
\end{equation}
The corresponding unitarily rescaled screw-ramp operator has kernel
$R_{\ell,L}(x,y)=wL(L|x-y|-\ell)_+$ and Hilbert--Schmidt norm
\begin{equation}\label{eq:ns55-ramp-norm}
 \|R_{\ell,L}\|_{\rm HS}
      =\frac{w}{\sqrt6}(L-\ell)_+^2.
\end{equation}
In particular, smoothing to the potential makes prime entry small
but does not make this term real analytic across $L=\ell$.
The actual bounded perturbation $\widetilde{\mathcal K}_L$ of
Proposition~\ref{prop:v118-log-dirichlet}, after this rescaling,
satisfies for each prime power $n_0$ with $\ell=\log n_0$
\begin{equation}\label{eq:ns55-full-k-jump}
 \liminf_{L\downarrow\ell,\,L>\ell}
 \|\widetilde{\mathcal K}_L-
             \widetilde{\mathcal K}_{\ell}\|_{2\to2}\ge w_{n_0}>0.
\end{equation}
Thus an analytic Fredholm argument in the raw bounded window
perturbation cannot be invoked at prime entry. No failure of
form or resolvent continuity is inferred.
\end{proposition}
\begin{proof}
For $\ell<L<2\ell$, put $b=\ell/L>1/2$. The shift interchanges
$L^2(0,1-b)$ and $L^2(b,1)$ by translation and is zero on the
middle interval. This proves its norm. For the potential kernel,
\[
 \|R_{\ell,L}\|_{\rm HS}^2
 =2w^2L^2\int_{\ell/L}^1(1-t)(Lt-\ell)^2\,dt
 =\frac{w^2}{6}(L-\ell)^4
\]
when $L>\ell$; it is zero otherwise. An operator-valued real-analytic
function vanishing for $L<\ell$ would also vanish just above $\ell$.

For the last assertion take normalized $u_L$ supported in the
right interval $(b,1)$ and project the output onto $(0,1-b)$.
The new shift has norm one on that input and output pair. For $L$
sufficiently close to $\ell$, no other prime power enters, and
all old shifts have lengths bounded strictly below one after
rescaling; their outputs from this shrinking right interval miss
the shrinking left interval, at both $L$ and $\ell$.
The scalar diagonal misses it too. The rescaled $k$ and pole
integral kernels are uniformly bounded near $L=\ell$, so their
restricted input--output norm is $O(1-b)$ by the Hilbert--Schmidt
bound. The projected difference therefore has norm at least
$w_{n_0}-O(1-b)$.
\end{proof}

\begin{proposition}[What fixed-window Fredholm reduction actually gives]
\label{prop:ns55-coupling-fredholm}
Fix $a$. Let $\mathcal A=\tfrac12L_\Delta^{\rm Dir}$ and
$A_a=\mathcal A+\mathcal K_a$ be the complete realization.
Choose $\beta$ so that $T=\mathcal A+\beta I\ge I$, and set
\[
 B=\mathcal K_a-\beta I,\qquad
 C=T^{-1/2}BT^{-1/2}.
\]
Then $C$ is compact and self-adjoint, and
\begin{equation}\label{eq:ns55-fredholm-equivalence}
 \ker A_a=\{0\}\quad\Longleftrightarrow\quad
 -1\notin\sigma(C).
\end{equation}
More generally the common-domain coupling family $T+zB$ is
invertible except at the discrete real set
\[
 \{-1/\nu:\nu\in\sigma(C)\setminus\{0\}\},
\]
and its inverse away from that set is
$T^{-1/2}(I+zC)^{-1}T^{-1/2}$. It is meromorphic in $z$,
with finite-rank principal parts. This does not exclude the actual
coupling $z=1$ from the exceptional set.
\end{proposition}
\begin{proof}
The reference operator has compact resolvent, so $T^{-1/2}$ is
compact and bounded. The asserted properties of $C$ follow.
If $(I+zC)v=0$ and $f=T^{-1/2}v$, then
$f=-zT^{-1}Bf\in\mathcal D(T)$ and $(T+zB)f=0$.
Conversely a nullvector gives $v=T^{1/2}f$ in the kernel of
$I+zC$. The same factorization gives the inverse when it exists.
The spectral theorem for the compact self-adjoint $C$ proves the
exceptional-set and meromorphy assertions. Its nonzero eigenvalues
have finite multiplicity and accumulate only at zero, so their
reciprocals have no finite accumulation.
\end{proof}

\paragraph{Named missing input.}
One still needs to exclude $-1$ from the spectrum of this specific
arithmetic compact operator at every required window, or prove
uniqueness in the shifted-primitive equation by a different mechanism.
The Fredholm reduction supplies neither estimate. Applying analytic
Fredholm theory directly to the compact potential operator
$\mathcal C_a$ from v1.54 would also be invalid: a compact operator
on an infinite-dimensional space is not Fredholm, whereas the
standard theorem concerns an identity plus a compact operator.
The preconditioning above repairs that issue at fixed $a$ but
leaves the physical exceptional point unresolved. These are
obstructions to proposed shortcuts, not to arithmetic injectivity.
